% !TEX root = ../main.tex
\chapter{Foreign Exchange}\label{ch:fx}
\begin{paracol}{2}

\begin{sources}
\begin{tabularx}{\linewidth}{@{}l l L@{}}
\toprule
\textbf{Segment} & \textbf{Length} & \textbf{Content}\\
\midrule
\seg{16}{1}{L16.1 FT: High frequency trading} & 4:41 & EBS curbs ``predatory'' high-frequency traders.\\
\seg{16}{2}{L16.2 Uncovered interest parity} & 2:22 & Three FX facts: CIP, UIP, EH; UIP and EH fail for the same reason.\\
\seg{16}{3}{L16.3 FX dealers under the gold standard, redux} & 5:03 & The two gold-standard diagrams and their exogenous bounds.\\
\seg{16}{4}{L16.4 Private FX dealing system} & 10:41 & The \$10 payment: matched-book and speculative dealers.\\
\seg{16}{5}{L16.5 Economics of the dealer function, speculative dealer} & 5:57 & Forward rate below the expected spot rate: expected profit.\\
\seg{16}{6}{L16.6 Economics of the dealer function, matched-book dealer} & 6:15 & $F/S$ rises with the book; through CIP, foreign term rates are bid up.\\
\seg{16}{7}{L16.7 Digression: Why do UIP and EH fail?} & 9:45 & Dealers must be paid: liquidity premia, time-varying with the pattern of payments.\\
\seg{16}{8}{L16.8 Central bank as FX dealer of last resort} & 16:09 & Central bank swaps; the deficit central bank as speculative dealer; when intervention helps.\\
\seg{16}{9}{L16.9 Reading: McCauley on internationalization of renminbi} & 10:29 & Three tracks of China's financial system; FX as hybrid of state and market.\\
\bottomrule
\end{tabularx}

\smallskip
\textbf{Files.} Transcripts \file{materials/transcripts/L16-P*.txt}; Mehrling's notes \mn{16}{1}.
\textbf{Reading.} R.~N.~McCauley, ``Renminbi internationalisation and China's financial development'', \emph{BIS Quarterly Review}, December 2011; Baba and Packer, BIS WP~267 (2008); Mehrling, ``Essential hybridity (2013) (both in ile{materials/readings/}).
\end{sources}

\section*{Motivation}
Under the gold standard the central bank ``put its arms around the system'': the gold points and Bank rate were explicit outside spreads (lecture~\ref{ch:globalliq}).
The modern system has no such explicit bounds; the challenge is to understand how it works \ts{16}{3}{4:12}. Mehrling calls this ``in a way the hardest
lecture of the whole semester'': more wheels are turning than in any other \ts{16}{9}{8:46}. The payoff: the failures of uncovered interest parity and of the
expectations hypothesis turn out to be the same fact \ts{16}{2}{1:25}.

% ------------------------------------------------------------------
\section{FT: high frequency trading}\label{sec:fx-ft}

\begin{casestudy}[FT, 2012: ``EBS hails success of HFT curbs'']
High frequency trading (``cyborg trading''): algorithms, often written by computer scientists and engineers, run on servers co-located next to the exchange
\ts{16}{1}{0:06}. HFT firms make money, but it is not clear they make markets: they take money from slower regular dealers, who may
protect themselves by widening the bid--ask spread, making markets less liquid \ts{16}{1}{1:09}. EBS, one of the largest currency trading platforms, took measures
against ``predatory practices'': the fifth decimal of a quote must end in 0 or 5, and limits on orders cancelled before execution. Volumes fell, but the chief
executive says the real traders, and liquidity, stayed \ts{16}{1}{2:11}.
\end{casestudy}

Why does it matter especially in FX? Dealers hedge a spot position by a forward position; if a human takes a second between the two trades, a computer can take
both first and sell the hedge back at a profit. Banks that discover this back off \ts{16}{1}{1:30}. Platforms protect their other
customers because otherwise those customers go elsewhere: an arms race, the ``war of the cyborgs versus the humans'', in a market of \$4 trillion a day
\ts{16}{1}{3:32}.

\begin{definition}[New concepts: high frequency trading]
Automated trading by algorithms that place and cancel orders in fractions of a second, exploiting speed (\emph{high frequency}: many trades per unit of time)
\ts{16}{1}{0:06}.
\end{definition}

% ------------------------------------------------------------------
\section{Three FX facts}\label{sec:fx-facts}

\begin{keyformulas}
With $S$ the spot and $F$ the forward exchange rate (FX per dollar), $R$ and $R^\ast$ the dollar and foreign term rates \ts{16}{2}{0:03}:
\begin{enumerate}
  \item \textbf{covered interest parity} (CIP), a strict arbitrage condition that more or less defines the forward rate, since a synthetic forward can be built
        from deposits:\quad $\dfrac{F}{S}=\dfrac{1+R^\ast}{1+R}$;
  \item \textbf{uncovered interest parity} (UIP): the forward rate equals the expected spot rate,\quad $F=\E[S_T]$; empirically false;
  \item \textbf{expectations hypothesis} (EH): the long rate equals the expected rollover of short rates,\quad $(1+R_{0,T}) = (1+R_{0,N})\,\E[1+R_{N,T}]$; empirically
        false.
\end{enumerate}
\end{keyformulas}

The claim of the lecture: taking a money view of foreign exchange explains why UIP and EH fail, and shows that they fail for exactly the same reason:
``these are not different facts, they're actually the same fact'' \ts{16}{2}{1:25}.

% ------------------------------------------------------------------
\section{The gold standard, redux}\label{sec:fx-redux}

The two diagrams of last lecture \ts{16}{3}{0:24}:
\begin{itemize}
  \item exchange rate (pounds per dollar) against the dealers' liquidity risk: as the dollar depreciates the private dealers' quotes move down until the gold
        point $x/y-\delta$, where the US central bank must defend the dollar; the rate fluctuates between $x/y-\delta$ and $x/y+\delta$ even though there is a mint
        par \ts{16}{3}{1:47};
  \item the rate of interest at the centre (the UK), an upward-sloping curve: stress moves the market rate until Bank rate, where the Bank of England backstops
        the system, subject to convertibility; under an external drain it must raise Bank rate or suspend \ts{16}{3}{2:27}.
\end{itemize}
These were exogenous bounds set by the central bank. The modern system does not have them, at least not explicitly \ts{16}{3}{4:12}.

% ------------------------------------------------------------------
\section{The private FX dealing system}\label{sec:fx-private}

Back to the balance sheets of lecture~\ref{ch:chartalism}. A surplus country has a \$10 ``due from'' item on a deficit country; neither is the US, but trade is
invoiced and settled in dollars, the world reserve currency, and the deficit country has no dollars \ts{16}{4}{0:24}.

\paragraph{Step 1: the payment.} An FX dealer buys $10S$ of foreign exchange (FX, the deficit country's domestic currency) from the deficit country and pays with a new
liability, \$10 spot, which the deficit country transfers to the surplus country \ts{16}{4}{1:29}.
Both countries are done; the dealer is not happy: short dollars, long FX, exposed to the spot rate \ts{16}{4}{3:18}.

\paragraph{Step 2: the hedge.} The spot dealer hedges in the forward market, which is like a matched pair of term deposits, short and long, in the two currencies:
it adds a $10S$ FX term deposit liability (at $R^\ast$) and a \$10 term deposit asset (at $R$). This is CIP at work
\ts{16}{4}{4:19}. If spot and forward rates move together it loses on one side what it gains on
the other; some basis risk remains \ts{16}{4}{6:02}. One payment required two trades: this is why FX trading is \$4 trillion a day \ts{16}{4}{6:22}.

\paragraph{Step 3: the counterparty.} Whoever is on the other side of the hedge, neither country, takes the opposite position: a \$10 term deposit liability at $R$
and a $10S$ FX term deposit asset at $R^\ast$ \ts{16}{4}{6:44}.

\begin{nfigure}
\begin{taccts}
\tacct[2.0cm]{Deficit country}{$-$FX $10S$}{$-$Due to \$10}\quad
\tacct[2.0cm]{Surplus country}{$-$Due from \$10\\$+$\$10 spot}{}\quad
\tacct[2.0cm]{Matched-book dealer}{FX $10S$ spot\\\$10 term ($R$)}{\$10 spot\\FX $10S$ term ($R^\ast$)}
\end{taccts}
\begin{taccts}
\tacct[2.0cm]{Speculative dealer}{FX $10S$ term ($R^\ast$)}{\$10 term ($R$)}
\end{taccts}
\caption{The private FX dealing system: one \$10 payment, two dealers \ts{16}{4}{1:50}.}\label{fig:fx-private}
\end{nfigure}

\begin{definition}[New concepts: matched-book and speculative FX dealer (recalled)]
The \emph{matched-book dealer} hedges a long FX spot position with a short FX forward position: not exactly matched, but hedged. The \emph{speculative dealer} is
naked, speculating in the forward market, not in spot: it does not need to come up with spot dollars on demand. Its position is equivalent to borrowing at the
dollar rate and lending at the foreign rate: a carry trade \ts{16}{4}{7:53}.
\end{definition}

This is what finance people are paid for: making mutually beneficial trade possible across currencies, and shaving off their share. One dealer may take all of
these on its books; or no dealer will, and a central bank does it \ts{16}{4}{9:35}.

% ------------------------------------------------------------------
\section{The economics of the dealer function}\label{sec:fx-dealers}

\paragraph{The speculative dealer} is exposed to price risk, the subject of Treynor's model \ts{16}{5}{0:03}. On the vertical axis put the \emph{value} of FX (the
inverse of the course's usual quote, so that going down means depreciation of the deficit country's currency) \ts{16}{5}{0:24}.
What matters is the forward rate locked in today relative to the expected spot rate at maturity: being long FX, the dealer wants to buy it cheap. If the forward
price of FX is below the expected spot, the difference is expected profit \ts{16}{5}{1:46}.
Prices fluctuate and the dealer may be wrong, but right more than half the time it makes money. The speculative dealer is paid to take the exchange risk off the
matched-book dealer's hands \ts{16}{5}{4:37}.

\paragraph{The matched-book dealer} has hedged price risk but holds liquidity risk: it borrows short and lends long in the world reserve currency, with dollar
liabilities that can be withdrawn \ts{16}{6}{0:03}. It expands its book only if it buys spot more cheaply than it sells forward:
in $F/S$ space it has an upward-sloping bid--ask curve, higher the larger the book; in a crisis the spread also widens a lot \ts{16}{6}{1:08}.

\begin{nfigure}
\begin{tikzpicture}[font=\scriptsize,x=0.82cm,y=1cm]
  \begin{scope}
  \draw[-{Stealth[length=4pt]}] (-0.2,0) -- (6.4,0) node[below left,align=right]{long FX forward\\position};
  \draw[-{Stealth[length=4pt]}] (0,-0.2) -- (0,4.4) node[above]{value of FX (\$ per FX)};
  \draw[dashed,thmcol] (0,3.2) -- (6.2,3.2) node[above left]{$\E[S_T]$ (expected spot)};
  \draw[very thick,defcol] (0,3.1) -- (6,1.1);
  \draw[very thick,intcol] (0,2.8) -- (6,0.8);
  \fill (4,1.77) circle (1.5pt);
  \draw[{Stealth[length=3pt]}-{Stealth[length=3pt]}] (4,1.8) -- (4,3.17);
  \node[right,align=left] at (4.05,2.5) {expected\\profit};
  \node[left] at (4,1.55) {$F$};
  \node[align=center] at (3,-1.0) {speculative dealer};
  \end{scope}
  \begin{scope}[xshift=8cm]
  \draw[-{Stealth[length=4pt]}] (-0.2,0) -- (6.4,0) node[below left,align=right]{size of book\\(liquidity exposure)};
  \draw[-{Stealth[length=4pt]}] (0,-0.2) -- (0,4.4) node[above]{$F/S$, i.e.\ $(1+R^\ast)/(1+R)$};
  \draw[very thick,defcol] (0,1.2) -- (6,3.4);
  \draw[very thick,intcol] (0,0.9) -- (6,3.1);
  \node[align=center] at (3,-1.0) {matched-book dealer};
  \node[align=center,font=\scriptsize\itshape] at (2.2,3.6) {outside spread: ?};
  \end{scope}
\end{tikzpicture}
\caption{The two private FX dealers \ts{16}{5}{3:35}\ts{16}{6}{1:30}. Unlike the gold standard, no outside spread is drawn: who sets it is the question
\ts{16}{6}{2:58}. Schematic.}\label{fig:fx-dealers}
\end{nfigure}

\begin{intuition}[Expected spot as the new mint par]
In the notes, the expected spot rate plays, under floating, the role the mint par ratio played under gold: deviations create dealer profit as long as there is some expectation of return
\mn{16}{3}. But under floating there is no long-run anchor; and if people believe UIP, they read the forward depreciation that arises to induce dealers to absorb imbalances as a forecast of
future spot depreciation: a source of volatility. ``If that forecast leaks into expectations of actual future spot, then the whole system can get unhinged'' \mn{16}{3}\mn{16}{7}.
\end{intuition}

\paragraph{Where is the outside spread?} Under gold the central bank set the bounds; here nobody seems to. The key is CIP: $F/S=(1+R^\ast)/(1+R)$. If $F/S$ is
driven up by the matched-book dealer seeking profit, the term rates must move too; holding the US rate $R$ constant, the foreign term rate $R^\ast$ is bid up, an
upward-sloping curve in interest-rate space \ts{16}{6}{3:39}. The deficit country's central bank cares about
those term rates. Central banks set overnight rates, not term rates; the conceptual link between the two is the expectations hypothesis: with no liquidity
premia, term rates are expected rollovers of overnight rates, so a credible move in the overnight rate moves term rates \ts{16}{6}{5:10}.

% ------------------------------------------------------------------
\section{Why do UIP and EH fail?}\label{sec:fx-fail}

\begin{intuition}[One fact, two failures]
\begin{itemize}
  \item \textbf{EH fails}: to absorb the payment imbalance, term rates are bid up above what EH predicts. That is what we observe: on average the term structure is
        not flat but upward sloping. Dealers must make a profit to make markets, and this is where one kind of dealer's profit comes from \ts{16}{7}{0:03};
  \item \textbf{UIP fails}: if the forward rate equalled the expected spot rate, no speculative dealer would take the risky trade with zero expected profit:
        ``uncovered interest parity must be false'', or dealers will not make the market \ts{16}{7}{1:05}.
\end{itemize}
Both failures come from the need to mobilize profit-making dealers of two kinds to handle international payment imbalances \ts{16}{7}{1:46}.
\end{intuition}

\paragraph{An inverted yield curve?} A student saw an inverted repo yield curve on a terminal. Inverted curves usually signal monetary tightening, the central
bank pulling up the short rate; in repo just then, probably collateral shortage, unrelated to this argument \ts{16}{7}{2:27}.
An inverted curve is a sign of severe money-market stress: people bid for spot money because they have payments to make. Central banks usually cap overnight rates
unless they want stress; overnight rates pull the one-week and one-month rates through term-structure arbitrage, while at long horizons something other than
liquidity sets prices \ts{16}{7}{3:49}.

\begin{finance}[When CIP itself broke down]
In the crisis the need to roll dollar funding was so strong that those who could not borrow Eurodollars borrowed yen and swapped them into dollars: the dollar
funding crisis became an FX crisis, and covered interest parity broke down, because under stress people stopped doing the arbitrage and looked after themselves.
The BIS documented it \ts{16}{7}{5:11}. The lecture's argument is about normal times \ts{16}{7}{5:52}.
\end{finance}

\begin{coursenote}
The BIS article is not named in class. Representative BIS studies of the 2007--08 dislocation of FX swap markets are N.~Baba, F.~Packer and T.~Nagano, ``The
spillover of money market turbulence to FX swap and cross-currency swap markets'', \emph{BIS Quarterly Review}, March 2008, and N.~Baba and F.~Packer,
``Interpreting deviations from covered interest parity during the financial market turmoil of 2007--08'', BIS Working Paper 267, 2008.

The second paper is on the course's reading list (\file{materials/readings/L16-BIS\_work267\_Baba\_Packer.pdf}). Its abstract makes two points:
\begin{itemize}
  \item The sharp and persistent deviations from CIP between dollar and euro were significantly associated with the difference in counterparty risk between European and US institutions.
  \item The ECB's dollar term auctions, funded by the Fed's swap lines, stabilised the FX swap market by lowering the volatility of those deviations.
\end{itemize}
\end{coursenote}

\paragraph{What risk premium?} Abstracting from liquidity, standard economics has no story for either failure; it can invoke ``time-varying risk premia'', but what is
the risk? Here it is clear: exposure to liquidity, funding, turnover risk \ts{16}{7}{6:13}. The premia vary over time not because
utility functions bounce around but because the pattern of payments does: with small imbalances the premia should nearly vanish, since dealers need no
incentive \ts{16}{7}{7:14}. Inefficiencies appear even in the most liquid market in the world, \$4 trillion a day: it is not about
volume but about the incentive for dealers \ts{16}{7}{7:54}.

\paragraph{Pushing the survival constraint forward.} The dealers move the deficit country's survival constraint on into the next day: ``I don't have any dollars today,
I'm going to have dollars tomorrow, I think'', and I pay dealers meanwhile. Since they do not know when the flow reverses, they must price as if they will hold to
maturity; ``that's where the prices come from'' \ts{16}{7}{8:36}.

% ------------------------------------------------------------------
\section{The central bank as FX dealer of last resort}\label{sec:fx-cb}

The central bank comes in when private dealers are done: prices have moved, and they will take no more risk even if prices move dramatically
\ts{16}{8}{0:03}. Add the two countries' central banks; for simplicity they deal with each other, not with a third country or the IMF
\ts{16}{8}{0:50}. The private agent in the deficit country still needs to sell FX for dollars \ts{16}{8}{2:12}:
\begin{enumerate}
  \item the deficit central bank borrows term dollars from the surplus central bank at a rate $R^\ast_G$ (``G'' for government, not necessarily a market rate),
        and the surplus central bank opens a \$10 spot deposit for it: a swap of IOUs, like all banking \ts{16}{8}{2:33};
  \item the deficit central bank sells the \$10 to its resident for $10S_G$ of FX, at a rate that need not be the market rate (``we like you, you are an important
        industry, or a connected fellow'') \ts{16}{8}{4:46};
  \item the resident pays the surplus country: the dollar deposit is transferred \ts{16}{8}{5:28}.
\end{enumerate}

\paragraph{Sterilization.} The surplus central bank's balance sheet, and its money supply, expand; if it dislikes that it sterilizes by selling a Treasury bill.
Net, it has sold an asset and lent the proceeds to another central bank: its government may be unhappy about lending to foreigners instead of holding Treasury
bills, which ``might ring a bell'' from the crisis \ts{16}{8}{6:14}. The deficit
central bank's sterilization is the mirror, an expansionary operation adding an FX term asset \ts{16}{8}{8:20}.

\begin{nfigure}
\begin{taccts}
\tacct[2.6cm]{Surplus central bank}{$+$\$10 term loan to deficit CB ($R^\ast_G$)\\$-$T-bill (sterilization)}{($+$\$10 spot deposit of deficit CB,\\then transferred; net zero)}\qquad
\tacct[2.6cm]{Deficit central bank (net)}{$+$FX $10S_G$ term}{$+$\$10 term borrowing ($R^\ast_G$)}
\end{taccts}
\caption{The central bank as FX dealer of last resort, net positions after sterilization \ts{16}{8}{7:18}.}\label{fig:fx-cb}
\end{nfigure}

\paragraph{The deficit central bank as speculative dealer.} After the cancellations the deficit central bank borrows term dollars and lends term FX: it is a
speculative dealer, with a naked forward position in its own currency, but without a profit motive; if there were expected profit, private dealers would be doing
it \ts{16}{8}{9:04}. It ends up as FX dealer of last resort as a consequence of its commitment to peg overnight
interest rates \ts{16}{8}{10:32}.

\begin{definition}[New concepts: FX dealer of last resort, sterilization]
\emph{FX dealer of last resort}: the central bank taking the naked forward position private speculative dealers will no longer take, as a backstop for the FX
market \ts{16}{8}{10:32}. \emph{Sterilization}: an offsetting operation (for example selling a Treasury bill) that cancels the effect of an
FX operation on the domestic money supply (from Latin \emph{sterilis}, ``barren'': the operation is made to produce no monetary effect) \ts{16}{8}{6:34}.
\end{definition}

\paragraph{When is intervention efficient?} Knowing it is the dealer of last resort, the central bank may intervene earlier: if private speculative dealers charge
``an arm and a leg'' for the naked exposure, pushing the forward rate away from its efficient level, the central bank can provide the forward hedge to the
matched-book dealers itself \ts{16}{8}{10:53}. Typically these are central banks of poor countries, or countries without
well-developed FX dealer markets, where premia can be very large: ``I know I'm a minor currency, but my people are just as important as your people''
\ts{16}{8}{11:54}. Prices are signals: of competitiveness, and, for term rates, of whether saving is scarce. But here it is liquidity, not
saving, that is scarce, and the central bank can make it less scarce, accepting zero profit, or even losses for a national purpose such as industrialization
\ts{16}{8}{12:36}.

\begin{coursenote}
Not a blanket endorsement of currency intervention. A central bank offering FX hedges to the whole world would cause bubbles and ``all kinds of problems'';
with small distortions, why bother? Deep, liquid private markets need to make money, and a central bank that replaces them must decide politically who gets
foreign exchange, or end up ``offering Goldman Sachs their next bonuses''. The aim is to facilitate the underlying trade \ts{16}{8}{14:19}.

The notes widen the options \mn{16}{5}\mn{16}{6}. Central banks that target overnight rates ``are inevitably drawn into serving as FX dealers of last resort'': first as speculative dealers selling
hedges to matched-book dealers, then as matched-book dealers themselves through swaps with other central banks. The surplus central bank lending to a foreign central bank rather than its own
government is a limiting case; other sources of reserves work too, regional liquidity pools or the IMF. And no other central bank is needed: the deficit central bank can serve as speculative dealer
for private hedgers, or offer forward cover directly to its citizens; in the limit, forward cover for all comers at the current spot rate amounts to fixing the exchange rate. The rates involved
($R^\ast_G$, $S_G$) are policy rates, reflecting non-commercial relationships. The danger of trading with all comers at non-commercial prices is arbitrage by speculators; the notes quote DeRosa
(2009) on currency crises, all preceded by large positions short dollars and long local currency that the whole market, ``not counting the central bank'', had to hedge at once. Hence: ``the
exchange rate is not a free variable left open for state choice, but neither is it a market price that always fully reflects fundamental valuation'' \mn{16}{7}.
\end{coursenote}

% ------------------------------------------------------------------
\section{McCauley on the internationalization of the renminbi}\label{sec:fx-rmb}

\begin{reading}[R.~N.~McCauley, Renminbi internationalisation and China's financial development (BIS, 2011)]
McCauley describes China's monetary, credit and bond markets as having three tracks \ts{16}{9}{0:03}:
\begin{enumerate}
  \item the \textbf{old track}, the controlled economy, still large: banks are told whom to lend to and at what rate \ts{16}{9}{0:24};
  \item \textbf{peripheral markets}: loosened at the edges, a money market, a small interbank market, bond and loan markets with market rates; some market
        signals are better than none, in a controlled transition \ts{16}{9}{1:11};
  \item a new \textbf{offshore track}, from the internationalization of the renminbi \ts{16}{9}{2:16}.
\end{enumerate}
\end{reading}

The analogy in McCauley's mind is the Eurodollar market of the late 1950s and early 1960s: deposits and loans in dollars offshore, outside US regulation, often a
way of evading capital controls and interest-rate ceilings (Regulation~Q) \ts{16}{9}{2:16}. Offshore and onshore
renminbi prices are not on top of each other, a sign of obstacles to arbitrage (as Eurodollar and Fed funds rates normally are, within ``gold points'' of a few
basis points); in stress, the third and fourth quarters of 2011, they move farther apart, as Eurodollar and Fed funds did in the crisis
\ts{16}{9}{3:38}. Offshore renminbi bonds yield less than onshore ones, even with lower credit quality, because
holders of offshore renminbi deposits can switch only into offshore bonds \ts{16}{9}{5:01}. The Chinese authorities seem to allow it, as they
allowed track two; but this is harder to turn off once developed \ts{16}{9}{5:41}.

The notes place the offshore track in Hong Kong, money and capital markets both; the first step will be breaking down the barriers between the tracks, and in the process the central bank's
role will shift ``from a gold-standard like price setting function to a modern price backstopping function'' \mn{16}{7}.

\paragraph{A reserve currency?} Behind the discussion is the renminbi as a world reserve currency. China is a surplus country: everyone has to make payments to it,
and the survival constraint is asymmetric, so its liabilities are money, as US liabilities were after the two world wars. Today the Eurodollar market has little to
do with the US trade balance \ts{16}{9}{6:03}.

\begin{intuition}[FX is inherently hybrid]
Foreign exchange is a negotiated relationship between the state and the market, different in different countries. The Fed basically lets the private market
work: there are deep spot and forward markets with the dollar on one side. Elsewhere there are no dealers and central banks make the markets; most countries are
in between. These are not ideological choices: near the top of the hierarchy you look like the US, near the bottom like the central-bank case, in the middle a
hybrid. ``It's where states meet states and markets meet markets'' \ts{16}{9}{7:04}. There is no
general rule that intervention or capital controls are always bad; prices need not be at their efficient levels. Economists of all political stripes abstract
from liquidity, ``but liquidity doesn't abstract from them'' \ts{16}{9}{9:08}.
\end{intuition}

\begin{reading}[Mehrling, ``Essential hybridity: a money view of FX'' (\emph{J.~Comparative Economics} 41, 2013)]
The written version of this lecture (\file{materials/readings/L16-Essential\_hybridity.pdf}). Economics treats the exchange rate as the relative price of goods, finance as the relative price of assets. The paper treats it instead as the \emph{relative price of moneys}, set in dealer markets by order flow.
\begin{itemize}
  \item \textbf{A spectrum of regimes.} At one extreme the central bank quotes buy and sell prices and absorbs the order flow on its own balance sheet. At the other extreme private dealers set prices. In the general, hybrid case private dealers make the market and the central bank acts as ``dealer of last resort''.
  \item \textbf{The central bank as speculative dealer.} The deficit country's central bank borrows term dollars, for instance from the Fed through a swap line, and sells them spot to its citizens. It ends up with a naked forward position, the same position a private speculative dealer would take for profit; it shoulders FX risk that no private speculator will take. In the limit it offers forward cover to all comers at the current spot rate, which amounts to fixing the exchange rate (pp.~360--361).
  \item \textbf{The exchange rate is ``essentially hybrid''.} It is neither a market price that always reflects fundamentals, nor a free policy instrument.
  \item \textbf{The hierarchy of currencies.} According to the BIS figures quoted, 84.9\% of FX volume has the dollar as one leg. Of \$4 trillion a day, \$1.5 trillion is spot and \$1.8 trillion FX swaps. ``The FX market is fundamentally a money market, not a capital market'' (p.~361).
  \item \textbf{Majors, minors and crosses.} The majors all have the dollar as one leg; the minors trade against a major; most minor cross pairs do not trade at all. Every currency is ultimately a promise to pay dollars, through a regional key currency (Fig.~5, p.~362).
\end{itemize}
\end{reading}

% ------------------------------------------------------------------
\flushside
\section*{Key formulas and ideas}
\begin{keyformulas}
\begin{tabularx}{\linewidth}{@{}l L@{}}
CIP & $F/S=(1+R^\ast)/(1+R)$: arbitrage, defines the forward rate (breaks down in crises)\\
UIP & $F=\E[S_T]$: fails, since speculative dealers need $F\neq\E[S_T]$ to earn expected profit\\
EH & long rate $=$ expected rollover of short rates: fails, since term rates are bid up to pay matched-book dealers\\
Private system & matched-book dealer hedges spot with forward; speculative dealer holds the naked forward (carry trade)\\
Outside spread & no explicit bound; via CIP the deficit country's term rate, linked by EH to the overnight rate the central bank sets\\
Last resort & deficit central bank borrows term dollars, lends term FX: a speculative dealer without profit motive\\
Hybridity & FX regime reflects place in the hierarchy: markets at the top, central banks at the bottom\\
\end{tabularx}
\end{keyformulas}

\section*{Exercises}

\begin{exercise}[The two dealers with numbers]
The spot rate is $S=1.25$ FX per dollar, the six-month rates are $R=1\%$ and $R^\ast=3\%$ (not annualized). Write the T-accounts of the four agents for the \$10
payment, compute $F$ from CIP, and the speculative dealer's profit in dollars if the spot rate in six months is 1.25, 1.27 or 1.30.
\end{exercise}

\begin{exercise}[UIP and expected profit]
Show that if $F=\E[S_T]$ the speculative dealer's expected profit is zero. If its required expected return is 0.5\% over six months, how far must $F$ be from
$\E[S_T]$? Translate this into a premium in $R^\ast$ via CIP.
\end{exercise}

\begin{exercise}[Central bank swap]
Redo the central bank case with numbers: the deficit central bank borrows \$10 at $R^\ast_G=2\%$ and sells it to a resident at $S_G=1.20$ when the market rate is
1.25. Write all T-accounts, including both sterilizations, and identify who bears the exchange risk and who receives the subsidy.
\end{exercise}

\begin{exercise}[Offshore and onshore]
Offshore renminbi bonds yield 3\% and comparable onshore bonds 4\%. Explain why arbitrage does not close the gap, what happens to the gap as controls are relaxed,
and compare with the Eurodollar--Fed funds spread before and during 2008.
\end{exercise}
